\RequirePackage{luatex85}
\documentclass[lualatex]{ltjsarticle}
\usepackage{luatexja} % 日本語対応のためのパッケージ
\usepackage{amsmath}
\usepackage{amssymb}
\usepackage{mathtools}
\usepackage{bm}
\usepackage[margin=1in]{geometry}

\title{
    RCD 空間から LCD 空間へ\\[0.5em]
    \large 低正則性理論が失った 20 年と、これからの 20 年
}

\author{Arnold Robinson}
\date{\today}

\begin{document}

\maketitle

\section{Introduction}

すべては 1904 年にポアンカレ（Poincar\'e）が提出した有名な問題に端を発する：

\begin{quote}
すべての単連結な閉 3 次元多様体は 3 次元球面と同相であるか？
\end{quote}

この問題は 20 世紀の幾何学の発展を根本から変え、数々の強力な道具や理論を生み出す原動力となった。ハミルトン（Hamilton）が 1980 年代に導入したリッチ流（Ricci flow）
\[
    \frac{\partial}{\partial t} g_t = -2\operatorname{Ric}(g_t)
\]
は、従来のトポロジー的手法とは全く異なる視点を提供した。すなわち、幾何構造自体を時間とともに進化させ、その進化を通じて局所的な複雑さを徐々に解消しようというアプローチである。リッチ流の思想は熱方程式と深いアナロジーを持つ。曲率は進化の過程で拡散し、幾何構造はより正則な状態へと収束していく。しかし、通常熱方程式と異なるのは、リッチ流が特異点を生み出す点である。

したがって、真に困難な問題はリッチ流の存在を証明することではなく、$t \nearrow T$ のときに多様体上で何が起こるのか、そして特異時刻の後にこの幾何学的対象をどのように捉え続けるのかを理解することであった。ペレルマン（Perelman）によるリッチ流の研究は、この計画を最終的に完結させた。エントロピー、非崩壊（non-collapsing）、サージェリー（surgery）、特異点解析といった一連の道具立てを通じて、3 次元多様体は制御された幾何構造を持つ部分へと分解可能であることが示され、こうしてポアンカレ予想およびより一般的なサーストンの幾何化予想が解決された。

しかし、今日の視点から見ると、リッチ流がもたらした最も重要な遺産はおそらくポアンカレ予想そのものではない。真に重要なのは次である：

\begin{quote}
我々は、滑らかな幾何学的対象が縮退した後に一体何に変化するのかについての系統的な研究を開始した。
\end{quote}

一連のリーマン多様体 $(M_i,g_i)$ が適切な曲率および幾何学的条件の下で収束するとき、個々の $M_i$ がすべて滑らかな多様体であったとしても、その極限 $(X,d)$ はもはや滑らかな多様体とは限らない。特異空間は人工的に作られた反例ではなく、自然な幾何学的進化や幾何学的極限において避けて通ることのできない対象なのである。ここに現代幾何解析における根本的なパラダイムシフトが起きた。すなわち、滑らかな空間の研究 $\longrightarrow$ 滑らかな空間の極限の研究、への移行である。

\section{リッチ流からリッチ極限空間へ}

リッチ曲率の下界が統一的に保たれるという条件下で、グロモフ（Gromov）のコンパクト性理論は距離空間のレベルでの予コンパクト性をもたらし、チェーガー＝コーディング（Cheeger--Colding）はこれらの極限空間の内部構造に関する系統的研究をさらに推し進めた。$\operatorname{Ric}_{g_i} \geq K g_i$ を仮定する。個々の $(M_i,g_i)$ が $n$ 次元滑らかなリーマン多様体であっても、そのグロモフ＝ハウスドルフ極限 $(X,d)$ には特異点が出現しうる。しかし、これらの特異空間は決して完全に任意のものではない。

チェーガー＝コーディング理論は、リッチ曲率の下界が極限空間の構造を極めて強い形で制限することを示している。正則集合（regular set）、特異集合（singular set）、接錐（tangent cones）、分解写像（splitting maps）、概分解定理（almost splitting theorem）、体積収束（volume convergence）などの理論が徐々に体系化され、完全なリッチ極限空間（Ricci limit space）理論が形成された。ここで非常に重要なことが起こった：

\begin{quote}
特異点それ自体が幾何学的対象となり始めたのである。
\end{quote}

これは、低正則性が構造の完全な喪失を意味するわけではないことを示している。それどころか、曲率条件そのものが極限空間の中に新たな正則性と剛性を生み出すのである。しかしながら、チェーガー＝コーディング理論は依然としてリーマン多様体の列を出発点としている。そこで、さらに大胆な疑問が自然と湧き上がってきた：

\begin{quote}
もし極限空間がリーマン多様体から由来するという事実を完全に忘れてしまったとしたら、距離と測度そのものから直接リッチ曲率を定義することはできないだろうか？
\end{quote}

\section{リッチ曲率から最適輸送へ}

ロット＝シュトゥルム＝ヴィラニ（Lott--Sturm--Villani）は、この問いに対して画期的な解答を与えた。彼らは、リッチ曲率の下界が確率測度空間上のワッサースタイン（Wasserstein）幾何を通じて特徴づけられることを見出した。測度距離空間 $(X,d,m)$ に対し、確率測度空間 $\mathcal{P}_2(X)$ とワッサースタイン距離 $W_2$ を考える。リッチ曲率の下界は、ワッサースタイン測地線に沿ったエントロピー汎関数の凸性に言い換えることができる。こうして、リッチ曲率 $\longrightarrow$ 最適輸送、という結びつきが生まれた。

リッチ曲率の概念は、2階微分構造に依存する局所的な対象から、距離、測度、確率分布の間の輸送構造によって記述される対象へと変貌を遂げた。ここから曲率＝次元条件（Curvature-Dimension condition）、すなわち $CD(K,N)$ が誕生した。この転換の重要性は、リッチ曲率に新たな定義を与えたという事実を遥かに超越している。それは次を意味していた：

\begin{quote}
本来は滑らかな微分幾何学に属していた構造が、微分構造に全く依存しない本質的な記述を持つ可能性がある。
\end{quote}

これが、真の意味での特異リッチ幾何学への扉を開いたのである。

\section{RCD 空間：曲率からエネルギーへ}

しかしながら、$CD(K,N)$ 空間は依然として広範すぎた。特に、$CD$ 条件だけではフィンズラー（Finsler）幾何学を排除することができない。したがって、真の意味でのリーマン幾何学の特異的拡張を得たいのであれば、ヒルベルト空間構造（Hilbertian structure）を特徴づける条件をさらに加える必要がある。ここに RCD 空間が誕生した。大雑把に言えば、次のように表される：
\[
    RCD(K,N) = CD(K,N) + \text{無限小ヒルベルト性（infinitesimal Hilbertianity）}.
\]
ここで無限小ヒルベルト性は、チェーガー・エネルギー（Cheeger energy）
\[
    \operatorname{Ch}(f) = \frac{1}{2} \int_X |Df|^2 \, dm
\]
の2次性を通じて表現される。このステップは特に重要である。

我々はリーマン計量を直接定義の中に放り込んだわけではない。むしろ、エネルギー構造がヒルベルト空間としての 2 次構造を持つことを要求し、それによって勾配、ラプラシアン（Laplace operator）、熱流（heat flow）などの解析的対象を再び舞台に呼び戻したのである。こうして非常に美しいルートが形成された：
\[
    \text{リーマン多様体} \longrightarrow \text{リッチ極限空間} \longrightarrow \text{RCD 空間}.
\]
かくして RCD 理論は、非常に成功を収めた合成リッチ幾何学（synthetic Ricci geometry）となった。それは大量のリッチ極限空間を含むだけでなく、適切な測度グロモフ＝ハウスドルフ収束の下で優れた安定性を持つ。熱流、チェーガー・エネルギー、ラプラシアン、ボクナー不等式（Bochner inequality）、ソボレフ解析、整流可能性（rectifiability）、分解定理、そして数々の古典的幾何解析の成果が、この枠組みの中でそれぞれの対応物を得た。ここまで、すべてが順調にいっているかに見えた。しかし、まさにこの場所から、より危険な問題が頭をもたげ始めるのである。

\section{失われた 20 年}

過去 20 年間の RCD 理論は、目覚ましい技術的成果を挙げた。問題は、RCD が成し遂げたことが不十分である点にあるのではない。問題はむしろ次にある：

\begin{quote}
RCD はあまりにも成功しすぎてしまった。
\end{quote}

それは極めて強力な道具立てを構築し、滑らかな構造を持たない空間上でも、勾配、ヘッセ行列（Hessian）、ラプラシアン、熱流、そして多くの古典的解析対象を再定義することを可能にした。その結果、研究は非常に安定したパターンを形成していった：古典的リーマン幾何学 $\longrightarrow$ 滑らかさの除去 $\longrightarrow$ 合成的定式化（synthetic formulation）の探索 $\longrightarrow$ 古典的定理の成立証明 $\longrightarrow$ さらなる解析学の構築 $\longrightarrow$ さらなる定理の証明。このサイクルは、美しく深遠な数学を数多く生み出した。しかし、それこそが、過去 20 年間の低正則性幾何学において最も警戒すべき問題なのだ。

\subsection{我々は一体、特異空間を研究しているのか、それともリーマン多様体を復元しているのか？}

RCD 空間は真の特異性を持ちうる。それは多様体ではなく、接錐は複雑な構造を示しうるし、局所次元は変動しうる。微分構造自体も弱い意味でしか存在しないかもしれない。しかしながら、多くの研究の根本的な目標は次を証明することであった：RCD 空間 $\approx$ ほとんど至る所でリーマン幾何（Riemannian geometry almost everywhere）。整流可能性理論はその最も典型的な例である。我々は座標近傍、ユークリッド的接錐、分解写像、そして弱微分構造を求め続け、最終的に空間がほとんど至る所（almost everywhere）の意味で良好なユークリッド的性質を持つことを証明する。もちろん、これは重要である。しかし、ここにはより根本的な問いが潜んでいる：

\begin{quote}
もし我々の最終的な目標が常に、「その特異空間は『ほとんど至る所でやはりリーマン的である』」ことを証明することなのだとしたら、我々は本当に特異性を研究しているのだろうか？
\end{quote}

我々が本来問うべき疑問は、おそらくその逆なのだろう：「特異性それ自体がどのような幾何学を宿しているのか？」、「特異集合はどれほど小さいか」ではなく。

\subsection{古典的定理の復元から新しい定理の創造へ}

過去 20 年間のもう一つの顕著な特徴は、古典的なリーマン幾何学の定理の多くが RCD 空間へと拡張されたことである。分解定理、体積錐が計量錐を意味すること（volume cone implies metric cone）、最大直径定理、オバタの定理（Obata theorem）、等周不等式、ソボレフ不等式、スペクトル収束などが、RCD の枠組みでそれぞれのバージョンを獲得した。これらの仕事が莫大な価値を持つことは言うまでもない。しかし、我々は 2 つの全く異なる問題を区別する必要がある：どのような古典的構造が極限の中でも保持されるのか？ そして、極限それ自体がどのような新しい構造を創造したのか？

前者の問題についてはすでに非常に豊かな研究がなされてきた。しかし後者の問題は、同程度の注目を浴びることからは程遠かった。これが奇妙な状況を生み出している：

\begin{quote}
我々は、RCD 空間がなぜ依然としてリーマン多様体に「似ている」のかについてはますます理解を深めているが、なぜそれらがもはやリーマン多様体「ではない」のかについては、同じようには明確に理解していない。
\end{quote}

言い換えれば、低正則性理論における「特異性」は、しばしば制御すべき誤差として扱われ、説明されるべき数学的現象としては扱われてこなかったのである。

\subsection{安定性は終着点となってしまったのか？}

RCD 理論の最も驚異的な性質の一つは、もちろんその安定性である。もし $(X_i,d_i,m_i) \longrightarrow (X,d,m)$ が適切な意味で収束し、各 $X_i$ が一様な RCD 条件を満たすならば、極限空間も対応する RCD 構造を保持する。これにより、極めて強力な閉性原理（closedness principle）が確立された：リッチ下界 + 幾何学的収束 $\Longrightarrow$ 極限におけるリッチ下界。しかし、安定性が答えているのは次である：古典的構造は縮退した後に何が残るのか？ それは次の問いには答えていない：縮退した後に何が生まれたのか？

この 2 つの問題を混同してはならない。もちろん、安定性は理論が備えていなければならない重要な性質である。しかし、ある理論が「構造が極限で保たれるかどうか」だけに囚われているとき、それは極限における新しい現象を、既存の構造の弱い極限として再解釈しがちである。ここに警戒すべき可能性が浮かび上がる：安定性は構造を保存する一方で、新規性を隠蔽してしまうかもしれない（stability may preserve structure while hiding novelty）。

\subsection{道具が問題を決定し始める}

RCD 理論は徐々に、非常に成熟した技術的閉回路を形成していった：最適輸送 $\longrightarrow$ エントロピー凸性 $\longrightarrow$ チェーガー・エネルギー $\longrightarrow$ 熱流 $\longrightarrow$ ボクナー不等式 $\longrightarrow$ 2階解析学。この理論は極めて美しい。しかし同時に、それは方法論的な強い引力を生み出した。ひとたびある問題が「エネルギー + 凸性 + 勾配流」の形に書き換えられるやいなや、RCD の機械装置（machinery）が自然に介入できるようになる。その結果、研究者はその機械装置で処理できるような問題を探し続ける傾向に陥りやすくなる。これは異常な現象を引き起こしかねない：

\begin{quote}
我々はますます強力な道具を手に入れているが、研究している問題そのものは、それらの道具で処理できるものにますます近づいている。
\end{quote}

最終的に、方法論が幾何学を決定し始める（methodology begins to determine the geometry）。これこそが、いわゆる「失われた 20 年」の最も警戒すべき本質かもしれない。

\section{真に困難な問題}

したがって、我々はより鋭い問いを改めて投げかけなければならない：

\begin{quote}
リーマン多様体、重み付きリーマン多様体、そしてリッチ極限空間の統一的拡張としての RCD 空間は、古典的なリーマン幾何学からは直接見えないような新しい幾何構造を一体生み出したのだろうか？
\end{quote}

ここで、RCD 空間に新しい現象がないと言っているわけではない。それどころか、それらは古典的な多様体には存在しない莫大な特異的挙動を備えている。問題は、それらの現象が通常、「正則性をいかに復元するか」、「特異性をいかに制御するか」、「安定性をいかに証明するか」、そして「古典的定理をいかに拡張するか」という枠組みに組織化されてしまう点にある。真に欠けているのは、別の研究方向である：特異性を消去するのではなく、特異性から出発すること。すなわち、「極限を生き延びたものは何か？」から、「極限によって何が創造されたのか？」へと転換すること。これこそがおそらく、過去 20 年間の RCD 理論において最も重要でありながら、最も正面から問われてこなかった問題なのだろう。

\section{なぜローレンツ幾何学は自然な実験場なのか}

もし我々が低正則性理論の生み出しうる真の新しい構造を探求したいのであれば、ローレンツ幾何学（Lorentzian geometry）は極めて自然な実験場である。リーマン幾何学もローレンツ幾何学も、ともに計量を核心に据えているが、両者の幾何学的本質は同一ではない。ローレンツ幾何学においては、我々は計量だけでなく、因果構造（causal structure）と時間分離（time separation）をもつ。時間分離を $\tau(x,y)$ と書こう。それは通常の距離関数ではない。それは因果関係 $x \leq y$ によって支配されている。したがって、ローレンツ幾何学は少なくとも同時に次を含んでいる：計量 + 因果構造 + 時間分離。

これは、RCD の思想をそのままローレンツの世界に持ち込むとき、最初から根本的な危険性が存在することを意味している：

\begin{quote}
我々はリッチ曲率の拡張には成功したかもしれないが、ローレンツ幾何学の拡張には本当の意味で成功していない可能性がある。
\end{quote}

もし我々が $CD(K,N)$ のローレンツ的アナロジーだけを研究するならば、得られるのは主に体積変化率、ヤコビアンの評価（Jacobian estimates）、そして最適輸送におけるエントロピー凸性である。これらの結果が重要であることは言うまでもない。しかし、それらが記述するのは主として体積の歪み（volume distortion）である。一方、ローレンツ幾何学の最も独特な構造は、因果性と時間分離に他ならない。もしある理論が $\tau(x,y)$、$J^+(x)$、$J^-(x)$、そして因果測地線を真に説明できないのであれば、我々はローレンツ幾何学の低正則性バージョンを構築しているのか、それとも単にローレンツの匂いをまとった CD 空間を作っているに過ぎないのか？ これは全く異なる 2 つの問題である。

\section{さらに危険な問題：偏微分方程式（PDE）}

問題は因果構造よりもさらに深いところにある。RCD 理論が巨大な成功を収めうる決定的な理由は、リーマン幾何学における PDE がエネルギー構造と極めて密接に結びついている点にある。滑らかなリーマン多様体上では、$\Delta f = \operatorname{div}(\nabla f)$ は楕円型作用素である。対応するボクナー恒等式は次のように書ける：
\[
\frac{1}{2}\Delta |\nabla f|^2 = |\operatorname{Hess}f|^2 + \langle \nabla f, \nabla\Delta f \rangle + \operatorname{Ric}(\nabla f,\nabla f).
\]
これは PDE + エネルギー + リッチ曲率を結びつけている。しかし、ローレンツ幾何学における対応するダランベール作用素（d'Alembert operator）$\Box_g$ は双曲型作用素である。したがって、$\Delta_g$ と $\Box_g$ の PDE 理論における性質には根本的な違いがある。

リーマン幾何学における基本的な解析構造は、楕円型 PDE、熱流、エネルギー最小化である。これに対し、ローレンツ幾何学におけるそれは、双曲型 PDE、波動伝播（wave propagation）、因果的進化（causal evolution）である。ゆえに：熱流 $\neq$ 波動伝播、勾配流 $\neq$ 因果的進化。もし我々が RCD における「チェーガー・エネルギー $\longrightarrow$ 熱流 $\longrightarrow$ ラプラシアン」という一連の構造をそのままローレンツ幾何学に移植するならば、ローレンツ幾何学の最も重要な部分を根本から見落としてしまう可能性が高い。ローレンツ幾何学における情報は、単に空間全体へ拡散していくわけではない。それは方向を持つ。情報は因果錐（causal cones）に沿って伝播する。したがって、将来の合成ローレンツ幾何学（synthetic Lorentzian geometry）は、計量測度構造と凸性の上だけに構築されることはおそらくできない。それは同時に次をエンコードできなければならない：因果性 + 時間分離 + 双曲型 PDE + リッチ曲率。

\section{アインシュタイン方程式：真の終着点}

ローレンツ幾何学には、リーマン幾何学が回避できないもう一つの問題がある。アインシュタイン方程式（Einstein field equation）である。滑らかなローレンツ多様体上では、
\[
    \operatorname{Ric}_{\mu\nu} - \frac{1}{2}R g_{\mu\nu} + \Lambda g_{\mu\nu} = 8\pi T_{\mu\nu}
\]
が成り立つ。ここでリッチ曲率とは、孤立した曲率条件ではない。それは力学方程式の一部なのだ。したがって、もし我々が真の合成時空理論（synthetic spacetime theory）を構築したいと願うのであれば、リッチ曲率の下界を定義するだけでは到底不十分である。我々が真に問わなければならないのは次である：

\begin{quote}
古典的な微分構造が存在しない状況下において、アインシュタイン方程式は依然として自然な意味を持つだろうか？
\end{quote}

これは RCD 理論における問題とは本質的に異なる。RCD が主に答えているのは、「リッチ曲率の下界がいかにして特異空間上で表現されるか」である。一方、ローレンツ幾何学が最終的に答えなければならないのは、「アインシュタインの力学がいかにして特異空間上で表現されるか」である。この 2 者間には巨大な隔たりが存在する。

\section{「定理の拡張」から「構造の発見」へ}

それゆえ、これからの 20 年間の問題は、単に「より多くのローレンツ幾何学の古典的定理を LCD 空間へいかに拡張するか」であるべきではない。むしろ、「低正則性が我々にどのような新しい構造を発見させることを強制するのか」であるべきなのだ。第 1 の研究アプローチは、古典的幾何学 $\longrightarrow$ 弱定式化 $\longrightarrow$ 合成的定理、というものである。これに対し、第 2 の研究アプローチは次のものである：古典的幾何学 $\longrightarrow$ 正則性の喪失 $\longrightarrow$ 新たな現象 $\longrightarrow$ 新たな構造。

この 2者の差異は、技術的なルートの違いではない。それらは実際には、全く異なる 2 つの数学的哲学を代表している。過去 20 年間、我々は「滑らかさが消え去った後、古典的リーマン幾何学はどれほど生き残るのか」という問いに首尾よく答えてきた。しかし、未来において真に答えるべき価値のある問いは、おそらくその逆なのだろう：「滑らかさが消え去った後、何が初めて出現するのか？」

\section{これからの 20 年の課題}

この観点から見ると、今後の低正則性ローレンツ幾何学には、少なくともいくつかの基本的な問題が存在している。

\subsection{リッチ曲率はいかにして再解釈されるべきか？}

もし底空間が微分構造を持たないならば、$\operatorname{Ric}$ はもはや各点で定義できるテンソルではない。エントロピーの凸性を介してのみリッチ曲率の下界を定義することは、アインシュタイン方程式に必要な情報のすべてをすでに含んでいるのだろうか？ それとも、$\operatorname{Ric}$ と因果幾何学を同時に見通すような新しい対象が存在すべきなのだろうか？

\subsection{時間分離は基本変数となるべきか？}

ローレンツ幾何学において、$\tau(x,y)$ は抽象的なローレンツ計量よりも真の基本変数に近いかもしれない。したがって、真の合成時空は、$(X,d,m)$ から $(X,\tau,m,\leq)$ のような何らかのものへと移行すべきなのかもしれない。これは、理論が計量測度幾何学から因果測度幾何学（causal measure geometry）へ向かうことを意味している。

\subsection{チェーガー・エネルギーのローレンツ的アナロジーは存在するのか？}

RCD の無限小ヒルベルト性は、最終的にチェーガー・エネルギーの 2 次性によって具現化された。それでは、ローレンツの世界における対応するエネルギー構造とは何だろうか？ ローレンツ計量自体が不定値の符号を持つため、自然な候補は必ずしも正定値のディリクレ形式（Dirichlet form）とは限らない。したがって、我々は全く異なるエネルギー構造、すなわち不定値エネルギー（indefinite energy）、あるいはハミルトン力学により近い構造を必要としているのかもしれない。

\subsection{アインシュタイン方程式は合成的方程式となりうるのか？}

最終的に最も重要な問題は、もしかすると次のものである：アインシュタイン方程式は滑らかさなしで生き残ることができるのか（Can Einstein's equation survive without smoothness?）？ もし答えが単に「合成リッチ曲率の下界を定義できる」というだけであれば、それはやはりリッチ幾何学の拡張に過ぎない。真にエキサイティングな答えは次のようなものであろう：アインシュタイン時空 $\longrightarrow$ 合成アインシュタイン時空（synthetic Einstein spacetime）。もしそれが可能になるならば、我々が手にするのはもはやローレンツ幾何学の弱いバージョンではなく、真に新しい幾何学理論かもしれないのだ。

\section{結語：20 年の継承から 20 年の創造へ}

ポアンカレ予想からリッチ流へ、リッチ流からリッチ極限空間へ、チェーガー＝コーディングからロット＝シュトゥルム＝ヴィラニへ、そして RCD 空間へ。現代幾何学は一世紀の歳月をかけて、滑らかな構造がない状況下でいかにリッチ曲率を議論し続けるかを学んできた。この道程は我々に教えている：滑らかさは幾何学の本質ではない、と。しかしそれは同時に、より困難な問いを残した：もし滑らかさが本質ではないのなら、何が本質なのだろうか？

過去 20 年間の RCD 空間の理論は、「古典的リーマン幾何学はどれだけ温存されうるか」という問いに極めて見事に答えてきた。しかし、我々はまだ真に答えてはいない：低正則性そのものは何を創造したのか？ もし過去の主要な任務が「RCD 空間 $\approx$ ほとんど至る所でリーマン幾何」を証明し続けることだったのなら、次の 20 年間で真に研究すべき問題はおそらくその逆である：RCD 空間 $\not\approx$ リーマン幾何。

我々が探求すべきなのは、「どこが依然としてユークリッド空間であるか」ではない。むしろ、「どこからユークリッド空間でなくなるのか」である。「どの古典的定理が依然として成立するか」ではなく、「どのような新しい定理が特異空間においてのみ出現しうるのか」である。「いかにして微分構造を復元するか」ではなく、「微分構造がなくなった後、どのような構造がそれに取って代わるのか」である。

これこそが、ローレンツ幾何学がこれほどまでに重要である理由なのだ。ローレンツの世界では、我々はリーマン理論のあらゆる道具をそのままもう一度築き上げることはできない。因果性は整流可能性によって掻き消されることはない。時間分離はチェーガー・エネルギーに取って代わられはしない。双曲型伝播は良好な熱流を持つという理由だけで消え去ることはない。アインシュタイン方程式も、合成リッチ曲率の下界が存在するというだけで勝手に出現したりはしない。

したがって、ローレンツ幾何学は、低正則性理論にその未完の任務を真に完遂することを強制するかもしれない：古典的構造の復元から、新しい構造の発見へ。もし過去 20 年間が、「滑らかさが消え去った後、リーマン幾何学はどれほど生き残れるか」に答えるものだったとするならば、これからの 20 年間はおそらく、「滑らかさが消え去った後、幾何学は一体何に変化するのか」に答え始めるべきなのだろう。

これこそがおそらく、低正則性幾何学における真にまだ始まっていない部分なのだ。

\end{document}
